\documentclass[11pt]{article}
\usepackage[margin=0.85in]{geometry}
\usepackage[T1]{fontenc}
\usepackage{amsmath,amssymb,booktabs,graphicx,xurl,hyperref}
\graphicspath{{outputs/task1/}{outputs/task2/}{outputs/task3/}{outputs/bonus/}}

\title{CS 498 HW2 Report\\Point Cloud Part Segmentation}
\author{Your Name (NetID)}
\date{Fall 2026}

\begin{document}
\maketitle

\section{Evaluation protocol}
State that training data fit model parameters, validation mean IoU selected
all settings and the best checkpoint, and test data were used only after all
choices were frozen. Confirm that you made one final test evaluation for
PointNet and one for the Transformer and did not tune after viewing them.

\section{Data, augmentation, and metrics}
Describe the supplied point representation and the augmentation you
implemented. Include \path{augmentation.png} and \path{class_frequency.png}.
State the transformation ranges you used and explain why the semantic labels
remain unchanged. Give your confusion-matrix convention, per-class IoU
equation, and mean IoU convention.

\section{PointNet baseline}
Using the fixed architecture from the handout, explain how the
\texttt{[B,N,6]} input becomes local \texttt{[B,256,N]} features, a global
\texttt{[B,256,1]} feature, and point-wise \texttt{[B,N,5]} logits. Include
\path{training_curves.png} and the Task 2
\path{predictions.png}. Report parameter count, best validation epoch, the
validation loss at that epoch, overall point accuracy, per-class IoU, and mean
IoU. Also report the aggregate test accuracy, per-class IoU, and mean IoU from
the final validation-selected checkpoint. Discuss
the training/validation loss gap and one failure case.

\section{Transformer}
Explain how query, key, and value projections let every point exchange
information with every other point. State the packed-QKV and per-head tensor
shapes for the fixed width-96, four-head architecture, and contrast this with
PointNet's single global maximum. Include the Task 3 \path{training_curves.png} and
\path{model_comparison.png}. The comparison must show at least two fixed
validation chairs, each with input RGB, ground truth, PointNet prediction, and
Transformer prediction.
Report the same metrics as for PointNet, together with parameter count and
training time, including the one final aggregate test evaluation. Compare the
curve shapes and both models under the shared data
and evaluation setup. Discuss which chair parts benefit most, the quadratic
cost of full self-attention, and one visible error.

\section{Limitations and improvements}
Discuss at least two limitations involving annotation quality, part imbalance,
point sampling, model capacity, or generalization. Propose one concrete next
experiment.

\section*{Optional Utonia linear-probe bonus}
If attempted, identify the official pretrained Utonia checkpoint and source,
describe the input preprocessing and frozen feature extraction, and specify
the point-wise linear head. Explain whether pretrained frozen features improve
over the two models trained from scratch. Use a library PCA routine such as
\texttt{torch.pca\_lowrank}; you do not need to implement PCA. Following the
DINOv2 feature-visualization convention, map the first three components to RGB.
Include \path{feature_pca.png}, the
linear-probe \path{training_curves.png}, and qualitative predictions for at
least two fixed validation chairs. Report best epoch, training time, trainable
parameter count, validation loss, overall point accuracy, per-class IoU, and
mean IoU. Discuss what structure is visible in the PCA colors, remembering
that PCA component signs are arbitrary. Submit your separate bonus code as
instructed in the final handout. This is the only section in which Gen-AI
leverage is allowed.

\section*{AI-use declaration}
Check one statement.

\begin{itemize}
  \item[$\square$] I did not use generative AI for this homework.
  \item[$\square$] I did not use generative AI for any required part. I used it
  only in the optional bonus marked ``Gen-AI leverage allowed''. Tool/model,
  purpose, and affected files: \underline{\hspace{1.4in}}
\end{itemize}

I confirm that generative AI did not write or complete my required code,
derive my required answers, generate my analysis, or write my required report.

\end{document}
